% !PRINTED 457
\weberpage{457}
Die gerade Linie $x_1 = 0$ schneidet die Curve in drei zusammenfallenden
Punkten bei $x_3 = 0$ und in einem vierten davon getrennten Punkte
$x_2 = 0$. Die Eckpunkte des Coordinatendreiecks sind also Wendepunkte
der Curve und die Seiten sind die Wendetangenten. Bezeichnen wir die
Seiten des Coordinatendreiecks mit 1, 2, 3 und die gegenüberliegenden
Ecken durch dieselben Ziffern, so ist die Seite 1 Wendetangente im
Punkt 2, die Seite 2 Wendetangente im Punkt 3 und die Seite 3
Wendetangente im Punkt 1.

Wir untersuchen nun die Lage der Pole und Axen in Bezug auf die
Grundcurve. Da die Eckpunkte des Coordinatendreiecks zu den
siebenzähligen Polen gehören, so schliessen wir zunächst,
\gesperrt{dass alle siebenzähligen Pole $P_7$ auf der Grundcurve
liegen.}

Es sind, wie die Eckpunkte selbst, alles \gesperrt{Wendepunkte} der Curve, und
diese ordnen sich in acht Wendepunktsdreiecke.

\gesperrt{Die dreizähligen Pole} $P_3$ \gesperrt{liegen gleichfalls auf der Grundcurve;
denn setzt man}
\[
x_2 = \varrho\,x_1, \quad x_3 = \varrho^2 x_1
\]
oder
\[
x_2 = \varrho^2 x_1, \quad x_3 = \varrho\,x_1,
\]
worin $\varrho$ eine cubische Einheitswurzel ist, so reducirt sich
$f(x_1, x_2, x_3)$ auf
\[
1 + \varrho + \varrho^2 = 0.
\]
Um die Schnittpunkte ihrer Verbindungslinie
\[
T = x_1 + x_2 + x_3 = 0
\]
mit der Grundcurve zu finden, setzt man $x_3 = -x_1 - x_2$, wodurch
$f(x_1, x_2, x_3) = 0$ in
\[
x_1^3(x_1+x_2) - x_2^3 x_1 + (x_1+x_2)^3 x_2 = (x_1^2+x_1x_2+x_2^2)^2 = 0
\]
übergeht. Da die linke Seite ein Quadrat ist, so fallen die vier
Schnittpunkte zweimal zu zweien zusammen, und es folgt, dass die
Linie $T$ eine Doppeltangente der Grundcurve ist.

\gesperrt{Die 28 Linien $T$ sind Doppeltangenten der Grundcurve; ihre
Berührungspunkte sind die 56 Pole $P_3$.}

Hieraus folgt beiläufig, dass die Punkte $P_7$ nicht auf den Linien
$T$ liegen, da eine Doppeltangente ausser den Berührungspunkten
keinen weiteren Schnittpunkt mit der Curve haben kann.

\gesperrt{Die sechs- und achtzähligen Pole liegen nicht auf der
Grundcurve.}
